\documentclass[11pt,a4paper]{article}
\usepackage{operationsform}
\begin{document}
\OPStitle{改善・振り返りチェックリスト}{TPL-CHK-045}{0.2.0}
\OPSsection{1}{確認項目}
\OPScheck{1}{最初に定義した目標と実測結果を同じ単位で比較した}
\OPScheck{2}{良かった点と困難を根拠つきの事実として記録した}
\OPScheck{3}{原因仮説、認識上の偏り、相反する見解を明示した}
\OPScheck{4}{改善項目を制度・技術・表現・運用に分類した}
\OPScheck{5}{実施の権限、費用と担当・期限を確認した}
\OPScheck{6}{改善確認の指標と次回の検証時点を定めた}
\OPScheck{7}{次の企画で参照可能な教訓を保存記録につないだ}
\OPSsection{2}{点検結果の記録}
\OPSblankshort{対象案件・対象版}
\OPSblankshort{担当者と点検日}
\OPSblanklong{未解決事項と対応}
\par\Needspace{4\baselineskip}\noindent\textbf{確認に用いた資料ID}\par\vspace{1mm}
{\fontsize{8.5}{13}\selectfont
\begin{tabularx}{\linewidth}{@{}XXXX@{}}
\textbf{対象・項目} & \textbf{内容・参照} & \textbf{担当・状態} & \textbf{備考}\\[2mm]
\strut & \strut & \strut & \strut\\[5mm]
\strut & \strut & \strut & \strut\\[5mm]
\end{tabularx}
}\vspace{2mm}\par
\end{document}
